\documentclass{amsart}
% For dealing with references we use the comment environment
\usepackage{verbatim}
\newenvironment{reference}{\comment}{\endcomment}
%\newenvironment{reference}{}{}
\newenvironment{slogan}{\comment}{\endcomment}
\newenvironment{history}{\comment}{\endcomment}

% For commutative diagrams we use Xy-pic
\usepackage[all]{xy}

% We use 2cell for 2-commutative diagrams.
\xyoption{2cell}
\UseAllTwocells

% We use multicol for the list of chapters between chapters
\usepackage{multicol}

% This is generally recommended for better output
\usepackage{lmodern}
\usepackage[T1]{fontenc}

% For cross-file-references

% Package for hypertext links:
\usepackage{hyperref}

% For any local file, say "hello.tex" you want to link to please

% Theorem environments.
%
\theoremstyle{plain}
\newtheorem{theorem}[subsection]{Theorem}
\newtheorem{proposition}[subsection]{Proposition}
\newtheorem{lemma}[subsection]{Lemma}

\theoremstyle{definition}
\newtheorem{definition}[subsection]{Definition}
\newtheorem{example}[subsection]{Example}
\newtheorem{exercise}[subsection]{Exercise}
\newtheorem{situation}[subsection]{Situation}

\theoremstyle{remark}
\newtheorem{remark}[subsection]{Remark}
\newtheorem{remarks}[subsection]{Remarks}

\numberwithin{equation}{subsection}

% Macros
%
\def\lim{\mathop{\mathrm{lim}}\nolimits}
\def\colim{\mathop{\mathrm{colim}}\nolimits}
\def\Spec{\mathop{\mathrm{Spec}}}
\def\Hom{\mathop{\mathrm{Hom}}\nolimits}
\def\Ext{\mathop{\mathrm{Ext}}\nolimits}
\def\SheafHom{\mathop{\mathcal{H}\!\mathit{om}}\nolimits}
\def\SheafExt{\mathop{\mathcal{E}\!\mathit{xt}}\nolimits}
\def\Sch{\mathit{Sch}}
\def\Mor{\mathop{\mathrm{Mor}}\nolimits}
\def\Ob{\mathop{\mathrm{Ob}}\nolimits}
\def\Sh{\mathop{\mathit{Sh}}\nolimits}
\def\NL{\mathop{N\!L}\nolimits}
\def\CH{\mathop{\mathrm{CH}}\nolimits}
\def\proetale{{pro\text{-}\acute{e}tale}}
\def\etale{{\acute{e}tale}}
\def\QCoh{\mathit{QCoh}}
\def\Ker{\mathop{\mathrm{Ker}}}
\def\Im{\mathop{\mathrm{Im}}}
\def\Coker{\mathop{\mathrm{Coker}}}
\def\Coim{\mathop{\mathrm{Coim}}}

% Boxtimes
%
\DeclareMathSymbol{\boxtimes}{\mathbin}{AMSa}{"02}

%
% Macros for moduli stacks/spaces
%
\def\QCohstack{\mathcal{QC}\!\mathit{oh}}
\def\Cohstack{\mathcal{C}\!\mathit{oh}}
\def\Spacesstack{\mathcal{S}\!\mathit{paces}}
\def\Quotfunctor{\mathrm{Quot}}
\def\Hilbfunctor{\mathrm{Hilb}}
\def\Curvesstack{\mathcal{C}\!\mathit{urves}}
\def\Polarizedstack{\mathcal{P}\!\mathit{olarized}}
\def\Complexesstack{\mathcal{C}\!\mathit{omplexes}}
% \Pic is the operator that assigns to X its picard group, usage \Pic(X)
% \Picardstack_{X/B} denotes the Picard stack of X over B
% \Picardfunctor_{X/B} denotes the Picard functor of X over B
\def\Pic{\mathop{\mathrm{Pic}}\nolimits}
\def\Picardstack{\mathcal{P}\!\mathit{ic}}
\def\Picardfunctor{\mathrm{Pic}}
\def\Deformationcategory{\mathcal{D}\!\mathit{ef}}

\setlength{\emergencystretch}{3em}
\begin{document}
\title[Stacks Project, Tag 09EE]{420 --- Wild inertia and the tame inertia character}
\maketitle
\noindent Copyright (C) 2005--2025 Johan de Jong. Stacks Project authors. Source: \href{https://stacks.math.columbia.edu/tag/09EE}{Tag 09EE}.

\noindent Editorial difficulty note (not author text). Difficulty H2: The author constructs a uniformizer-independent character of inertia into roots of unity and identifies its kernel through the associated graded ring. A filtration calculation forces kernel elements to have p-power order, and a formally etale residue-field splitting proves the second-order description under separability..

\noindent Editorial provenance note (not author text). History: excerpt from revision \texttt{a0444\allowbreak 6e57e\allowbreak c1fbc\allowbreak 252a8\allowbreak 71afc\allowbreak ec775\allowbreak 2fb28\allowbreak 07b14}, more-algebra.tex, lines 33250--33406. Reference presentation was adapted; original-author context and core support blocks were retained or added. The mathematical wording is unchanged. Licensed under GNU FDL 1.2 or later; no invariant sections or cover texts. See ../licenses/COPYING. Context blocks reproduce author definitions and supporting results; they are not additional counted problems.

\begin{definition}
\label{fields-definition-insep-degree}
Let $E/F$ be an algebraic field extension. Let $E_{sep}$ be the subextension
found in Lemma \href{https://stacks.math.columbia.edu/tag/030K}{lemma separable first (Tag 030K)}.
\begin{enumerate}
\item The integer $[E_{sep} : F]$ is called the {\it separable
degree} of the extension. Notation $[E : F]_s$.
\item The integer $[E : E_{sep}]$ is called the {\it inseparable
degree}, or the {\it degree of inseparability} of the extension.
Notation $[E : F]_i$.
\end{enumerate}
\end{definition}

\begin{lemma}
\label{lemma-galois-inertia}
Let $A$ be a discrete valuation ring with fraction field $K$.
Let $L/K$ be a finite Galois extension with Galois group $G$.
Let $B$ be the integral closure of $A$ in $L$. Let $\mathfrak m \subset B$
be a maximal ideal. The inertia group $I$ of $\mathfrak m$
sits in a canonical exact sequence
$$
1 \to P \to I \to I_t \to 1
$$
such that
\begin{enumerate}
\item if $D$ is the decomposition group we have
$$
P = \{\sigma \in I \mid \sigma|_{\mathfrak m/\mathfrak m^2} =
\text{id}_{\mathfrak m/\mathfrak m^2}\}
=
\{\sigma \in D \mid \sigma \text{ acts trivially on }\text{Gr}_\mathfrak m(B)\}
$$
\item $P$ is a normal subgroup of $D$,
\item $P$ is a $p$-group if the characteristic of $\kappa_A$ is
$p > 0$ and $P = \{1\}$ if the characteristic of $\kappa_A$ is zero,
\item $I_t$ is cyclic of order the prime to $p$ part of the integer $e$,
\item there is a canonical isomorphism
$\theta : I_t \to \mu_e(\kappa(\mathfrak m))$, and
\item
$P =
\{\sigma \in D \mid \sigma|_{B/\mathfrak m^2} = \text{id}_{B/\mathfrak m^2}\}$
if $\kappa(m)$ is separable over the residue field of $A$.
\end{enumerate}
Here $e$ is the integer of Lemma \href{https://stacks.math.columbia.edu/tag/09EB}{lemma galois conclusion (Tag 09EB)}.
\end{lemma}

\begin{proof}
Recall that $|G| = [L : K] = nef$, see Lemma \href{https://stacks.math.columbia.edu/tag/09EB}{lemma galois conclusion (Tag 09EB)}.
Since $G$ acts transitively on the set
$\{\mathfrak m_1, \ldots, \mathfrak m_n\}$ of maximal ideals of $B$
(Lemma \href{https://stacks.math.columbia.edu/tag/09EA}{lemma galois (Tag 09EA)})
and since $D$ is the stabilizer of an element we see that $|D| = ef$.
By Lemma \href{https://stacks.math.columbia.edu/tag/09ED}{lemma galois galois (Tag 09ED)} we have
$$
ef = |D| = |I| \cdot |\text{Aut}(\kappa(\mathfrak m)/\kappa)|
$$
where $\kappa$ is the residue field of $A$.
As $\kappa(\mathfrak m)$ is normal over $\kappa$ the order of
$\text{Aut}(\kappa(\mathfrak m)/\kappa)$ differs from $f$ by
a power of $p$ (see
Fields, Lemma \href{https://stacks.math.columbia.edu/tag/09HS}{lemma normal and automorphisms (Tag 09HS)}
and discussion following
Fields, Definition \ref{fields-definition-insep-degree}).
Hence the prime to $p$ part
of $|I|$ is equal to the prime to $p$ part of $e$.

\medskip\noindent
Set $C = B_\mathfrak m$. Then $I$ acts on $C$ over $A$ and trivially
on the residue field of $C$. Let $\pi_A \in A$ and $\pi_C \in C$ be
uniformizers. Write $\pi_A = u \pi_C^e$ for some unit $u$ in $C$.
For $\sigma \in I$ write $\sigma(\pi_C) = \theta_\sigma \pi_C$ for some
unit $\theta_\sigma$ in $C$. Then we have
$$
\pi_A = \sigma(\pi_A) = \sigma(u) (\theta_\sigma \pi_C)^e
= \sigma(u) \theta_\sigma^e \pi_C^e = \frac{\sigma(u)}{u} \theta_\sigma^e \pi_A
$$
Since $\sigma(u) \equiv u \bmod \mathfrak m_C$ as $\sigma \in I$
we see that the image $\overline{\theta}_\sigma$
of $\theta_\sigma$ in $\kappa_C = \kappa(\mathfrak m)$
is an $e$th root of unity.
We obtain a map
\begin{equation}
\label{equation-inertia-character}
\theta : I \longrightarrow \mu_e(\kappa(\mathfrak m)),\quad
\sigma \mapsto \overline{\theta}_\sigma
\end{equation}
We claim that $\theta$ is a homomorphism of groups and independent
of the choice of uniformizer $\pi_C$. Namely, if $\tau$ is a second
element of $I$, then
$\tau(\sigma(\pi_C)) = \tau(\theta_\sigma \pi_C) =
\tau(\theta_\sigma) \theta_\tau \pi_C$, hence
$\theta_{\tau \sigma} = \tau(\theta_\sigma) \theta_\tau$ and
since $\tau \in I$ we conclude that
$\overline{\theta}_{\tau \sigma} =
\overline{\theta}_\sigma \overline{\theta}_\tau$.
If $\pi'_C$ is a second uniformizer, then we see
that $\pi'_C = w \pi_C$ for some unit $w$ of $C$ and
$\sigma(\pi'_C) = w^{-1}\sigma(w)\theta_\sigma \pi'_C$,
hence $\theta'_\sigma = w^{-1}\sigma(w)\theta_\sigma$,
hence $\theta'_\sigma$ and $\theta_\sigma$
map to the same element of the residue field as before.

\medskip\noindent
Since $\kappa(\mathfrak m)$ has characteristic $p$, the group
$\mu_e(\kappa(\mathfrak m))$ is cyclic of order at most the prime
to $p$ part of $e$ (see Fields, Section \href{https://stacks.math.columbia.edu/tag/09HW}{section roots of 1 (Tag 09HW)}).

\medskip\noindent
Let $P = \Ker(\theta)$. By construction the elements of $P$ are exactly
the elements of $I$ which act trivially on
$\mathfrak m/\mathfrak m^2 = \text{Gr}_\mathfrak m^1(B)$,
i.e., $P = \{\sigma \in I \mid \sigma|_{\mathfrak m/\mathfrak m^2} =
\text{id}_{\mathfrak m/\mathfrak m^2}\}$. Also
$I$ consists of the elements of $D$ which act trivially on
$\kappa(\mathfrak m) = B/\mathfrak m$.
Since the graded ring $\text{Gr}_\mathfrak m(B)$ is generated
by $\text{Gr}^1_\mathfrak m(B) = \mathfrak m/\mathfrak m^2$ over
$\text{Gr}^0_\mathfrak m(B) = B/\mathfrak m$, we conclude that
$P = \{\sigma \in D \mid \sigma \text{ acts trivially on }
\text{Gr}_\mathfrak m(B)\}$.
Thus (1) is true. This implies (2) as $P$ is the kernel
of the homomorphism $D \to \text{Aut}(\text{Gr}_\mathfrak m(B))$.
If we can prove (3), then parts (4) and (5) will follow as $I_t$
will be isomorphic to $\mu_e(\kappa(\mathfrak m))$ as the arguments above show
that $|I_t| \geq |\mu_e(\kappa(\mathfrak m))|$.

\medskip\noindent
Thus it suffices to prove that the
kernel $P$ of $\theta$ is a $p$-group.
Let $\sigma$ be a nontrivial element of the kernel. Then $\sigma - \text{id}$
sends $\mathfrak m_C^i$ into $\mathfrak m_C^{i + 1}$
for all $i$. Let $m$ be the order of $\sigma$. Pick $c \in C$ such
that $\sigma(c) \not = c$. Then $\sigma(c) - c \in \mathfrak m_C^i$,
$\sigma(c) - c \not \in \mathfrak m_C^{i + 1}$ for some $i$ and
we have
\begin{align*}
0
& =
\sigma^m(c) - c \\
& =
\sigma^m(c) - \sigma^{m - 1}(c) + \ldots + \sigma(c) - c \\
& =
\sum\nolimits_{j = 0, \ldots, m - 1} \sigma^j(\sigma(c) - c) \\
& \equiv
m(\sigma(c) - c) \bmod \mathfrak m_C^{i + 1}
\end{align*}
It follows that $p | m$ (or $m = 0$ if $p = 1$). Thus every element of the
kernel of $\theta$ has order divisible by $p$, i.e., $\Ker(\theta)$
is a $p$-group.

\medskip\noindent
Proof of (6). Assume $\kappa(\mathfrak m)/\kappa$ is separable.
In this case $f = |D/I|$ by Lemma \href{https://stacks.math.columbia.edu/tag/09ED}{lemma galois galois (Tag 09ED)}
and $I$ has order $e$. If $e = 1$, then $P = I = \{\text{id}\}$
and the result is true. If $e > 1$, then
$B/\mathfrak m^2 = C/\pi_C^2C$ is a $\kappa$-algebra
and a small extension of $\kappa(\mathfrak m)$.
Because $\kappa(\mathfrak m)$ is formally \'etale over $\kappa$
(Algebra, Lemma
\href{https://stacks.math.columbia.edu/tag/090W}{lemma characterize separable algebraic field extensions (Tag 090W)})
there is a unique $\kappa$-algebra map
$\kappa(\mathfrak m) \to B/\mathfrak m^2$
right inverse to $B/\mathfrak m^2 \to \kappa(\mathfrak m)$.
In other words, there is a $D$-equivariant isomorphism
$B/\mathfrak m^2 = \mathfrak m/\mathfrak m^2 \oplus \kappa(\mathfrak m)$
of $\kappa$-algebras. This immediately shows that
$P =
\{\sigma \in D \mid \sigma|_{B/\mathfrak m^2} = \text{id}_{B/\mathfrak m^2}\}$
in this case.
\end{proof}
\end{document}
